\documentclass[11pt,a4paper]{article}

\usepackage[utf8]{inputenc}
\usepackage[T1]{fontenc}
\usepackage[french]{babel}
\usepackage{amsmath,amssymb}
\usepackage[margin=2.3cm]{geometry}
\usepackage{booktabs}
\usepackage{enumitem}

\title{Solution d'optimisation du planning médical}
\author{}
\date{}

\begin{document}
\maketitle

\section{Ensembles et indices}

\begin{itemize}[leftmargin=*]
    \item $M$ : ensemble des médecins ;
    \item $A$ : ensemble des activités médicales ;
    \item $J = \{\text{Lundi},\ldots,\text{Dimanche}\}$ : jours de la semaine ;
    \item $C = \{\mathrm{AM},\mathrm{PM},\mathrm{GA}\}$ : créneaux horaires.
\end{itemize}

Les activités considérées sont : pédiatrie, urgences, consultations,
exploration, astreinte, maternité et néonatologie.

\section{Paramètres}

\begin{itemize}[leftmargin=*]
    \item $b_{a,j,c}$ : nombre de médecins nécessaires pour l'activité $a$,
    le jour $j$ et le créneau $c$ ;
    \item $q_{m,a} \in \{0,1\}$ : vaut $1$ si le médecin $m$ possède la
    compétence nécessaire pour l'activité $a$ ;
    \item $v_{m,j,c} \in \{0,1\}$ : vaut $1$ si le médecin $m$ est absent
    le jour $j$ au créneau $c$ ;
    \item $P$ : activités prioritaires, soit urgences, néonatologie,
    pédiatrie, astreinte et maternité.
\end{itemize}

\section{Variables de décision}

\begin{align*}
    x_{m,a,j,c} &\in \{0,1\} &&\text{: médecin $m$ affecté à $a$ au créneau $(j,c)$},\\
    r_{a,j,c} &\in \mathbb{Z}_{+} &&\text{: nombre de remplaçants utilisés},\\
    h_m &\in \mathbb{Z}_{+} &&\text{: nombre de créneaux travaillés par $m$},\\
    d_m &\in \mathbb{R}_{+} &&\text{: écart absolu de charge du médecin $m$ à la moyenne}.
\end{align*}

\section{Fonction objectif}

La fonction objectif minimise d'abord le nombre total de remplaçants. Une
pénalité supplémentaire est appliquée aux remplacements des activités
prioritaires. Enfin, l'écart entre la charge maximale et la charge minimale
est réduit afin d'équilibrer le planning :

\begin{align*}
\min Z ={}& 100000 \sum_{a\in A}\sum_{j\in J}\sum_{c\in C} r_{a,j,c} \\
&+ \sum_{a\in A}\sum_{j\in J}\sum_{c\in C}
    w_a r_{a,j,c} \\
&+ 10 \sum_{m\in M}d_m,
\end{align*}

où

\[
w_a =
\begin{cases}
1000, & \text{si } a\in P,\\
100, & \text{sinon.}
\end{cases}
\]

Le coefficient $100000$ rend la minimisation du nombre de remplaçants
prioritaire par rapport aux autres critères.

\section{Contraintes}

\subsection{Couverture exacte des besoins}

Chaque besoin doit être couvert exactement par des médecins titulaires ou des
remplaçants :

\[
\sum_{m\in M} x_{m,a,j,c} + r_{a,j,c} = b_{a,j,c}
\qquad \forall a\in A,\;j\in J,\;c\in C.
\]

Lorsqu'aucun médecin n'est requis, aucune variable de remplacement n'est
créée pour le créneau concerné.

\subsection{Compétences médicales}

Un médecin ne peut être affecté qu'à une activité pour laquelle il est
compétent :

\[
x_{m,a,j,c} \leq q_{m,a}
\qquad \forall m\in M,\;a\in A,\;j\in J,\;c\in C.
\]

Dans l'implémentation, les variables correspondant à une compétence absente
ne sont pas créées.

\subsection{Absences et vacances}

Un médecin absent ne peut recevoir aucune affectation :

\[
x_{m,a,j,c} \leq 1-v_{m,j,c}
\qquad \forall m\in M,\;a\in A,\;j\in J,\;c\in C.
\]

\subsection{Charge hebdomadaire}

La charge de chaque médecin est définie par :

\[
h_m = \sum_{a\in A}\sum_{j\in J}\sum_{c\in C}x_{m,a,j,c}
\qquad \forall m\in M.
\]

La limite hebdomadaire est :

\[
h_m \leq 10
\qquad \forall m\in M.
\]

\subsection{Un seul créneau par jour}

Chaque médecin ne peut travailler qu'un seul créneau dans une même journée :

\[
\sum_{a\in A}\sum_{c\in C}x_{m,a,j,c} \leq 1
\qquad \forall m\in M,\;j\in J.
\]

\subsection{Repos après une garde}

Une garde de nuit interdit toute affectation pendant les trois créneaux du
jour suivant. Le planning étant hebdomadaire, le jour suivant dimanche est
le lundi :

\[
\sum_{a\in A}x_{m,a,j,\mathrm{GA}}
+ \sum_{a\in A}\sum_{c\in C}x_{m,a,j^+,c}
\leq 1
\qquad \forall m\in M,\;j\in J,
\]

où $j^+$ désigne le jour suivant avec bouclage de dimanche vers lundi.

\subsection{Équilibrage des charges}

On définit la charge moyenne par :

\[
\overline{h} = \frac{1}{|M|}\sum_{m\in M}h_m.
\]

Les écarts absolus à la moyenne sont linéarisés par :

\[
h_m-\overline{h} \leq d_m,\qquad
\overline{h}-h_m \leq d_m
\qquad \forall m\in M.
\]

La fonction objectif minimise alors la somme des écarts absolus
$\sum_{m\in M}d_m$.

\section{Résultat de résolution}

Le modèle a été résolu avec PuLP et le solveur CBC.

\begin{center}
\begin{tabular}{lr}
\toprule
\textbf{Indicateur} & \textbf{Valeur}\\
\midrule
Statut & Optimal\\
Valeur de la fonction objectif & 2\,419\,675{,}56\\
Nombre total de remplaçants & 24\\
Charge moyenne & 4,89 créneaux\\
Charge maximale observée & 7 créneaux\\
Charge minimale observée & 0 créneau\\
\bottomrule
\end{tabular}
\end{center}

\end{document}